\section{Discussion}\label{sec:discussion}

\subsection{Why a reader would use this, and on what evidence}

The question a reader brings is why to use this rather than what they already
have. The answer is in three parts and each has a different kind of support
behind it.

\emph{Because the quantity you must pin down is not the one you think.} The
interval depends on the two scales only through their ratio, and a requirement
stated on the variance component overstates what the band needs by $\rho^{-4}$ in
the number of populations. This is algebra, verified numerically, and it applies
whatever construction is used. It is also, by \eqref{eq:necessary}, a property of
the problem: the same $\rho^4/\{\varepsilon^2(1-\rho^2)^2\}$ appears in a
two-point lower bound on \emph{any} procedure that certifies the ratio, within a
factor of about $2.5$ of what our estimator achieves.

\emph{Because how an interval reacts to scale uncertainty is a property of its
form, and forms differ.} Proposition~\ref{prop:exchange} is exact, and it says
the deconvolution form pays $\lambda^{\rho^2/2}$ where every construction
proportional to a square-root scale bound pays $\lambda^{1/2}$. That covers a
class---normal-theory intervals, EBLUP prediction intervals for a target with a
new random effect, parametric bootstrap intervals built on them---rather than one
comparator. \textbf{It does not rank them absolutely}: the oracle width ratio
differs between constructions and the elasticity does not supply it. What it
supplies is the part that reverses, and the reversals are what a practitioner
would otherwise experience as unpredictability.

\emph{Because the resulting procedure is computable and its cost is measured.}
Under Assumption~\ref{as:R} it has the guarantee~\eqref{eq:guarantee}, the
supremum is bounded rather than searched, and against a comparator sharing the
target, information and level it is 14 to 22\% narrower in simulation and
narrower in 44 of 45 real survey configurations.

\subsection{What is not established}

\paragraph{The gap between the theorem and the application is the paper's main
weakness, and naming it does not close it.} Assumption~\ref{as:R} is a normal
model with a known centre and independent populations. In the European Social
Survey the centre is estimated, the scores carry a measured dependence of about
$-1/(n_r-1)$, the standardised deviations are badly non-normal for two of three
pre-specified items, and the $\chi^2$ law for a design-variance estimator is not
established for complex designs \citep{mcgovern2026}. We report the ESS analysis
as an application under those approximations with the guarantee not attached, and
a reader who needs a guarantee on survey data does not get one here.

\paragraph{Given the same structure information, a component-wise construction is
close.} Against a component-wise bound that estimates the same structure and
spends the same budget, the ratio construction is 0.7 to 2.3\% narrower. Most of
the gain over a structure-free baseline is the use of a variance structure, not
the ratio formulation. We prefer the ratio form because it has one fewer budget
piece, a shorter chain of sufficient conditions, and---among the constructions we
implemented---a conservative computational bound over a continuous structure
confidence set. That no such bound has been built for the component-wise form is
not evidence that none exists.

\paragraph{We do not benchmark against the conformal construction that avoids
the problem.} \citet{zhang2026conformal} obtain conditional validity for
in-sample area expectations without estimating the sampling variances, and their
reported half-widths grow with the design share where a direct interval's do not.
We have not run their construction against ours, because the two cover different
objects and a width comparison across targets is not informative. A reader whose
target is in-sample should read that comparison as unmade, not as decided here.

\paragraph{The lower bound is about certifying the ratio, not about width.}
Equation~\eqref{eq:necessary} says no procedure can pin $\kappa$ down with fewer
degrees of freedom. It does \emph{not} say no interval can be narrower: an
interval may be valid without certifying $\kappa$, by being conservative in some
other way, and we have not proved a width lower bound. Nothing here establishes
that the proposed interval is optimal. In the equal-variance submodel the
inferential problem is the classical one for an intraclass correlation, where
both the exact interval and the design trade-off are known; what the bound adds
is the assurance that \eqref{eq:twoaxis} is not an artefact of the estimator we
happened to write down.

\paragraph{The gain is conditional on the structure family.} A confidence set for
a structure parameter is not protection against the family being wrong. Where the
family fails, containment fails while the final coverage does not, and the latter
is downstream slack rather than a guarantee.

\paragraph{And what the registered test refuted.} Equation~\eqref{eq:S4}'s
invariance in the design share is asymptotic, not exact, and its constant needed
a $1+1.9/K$ correction whose value we fit rather than derive. The order-statistic
expansion would supply it; we have not done that.

\subsection{Scope}

The target here is a scalar population summary. We do not place a
whole-distribution simultaneous band and the scalar result under one guarantee:
extending Proposition~\ref{prop:cert} to coordinate-varying scales requires a
simultaneous bound on a vector of correction factors, with a multiplicity cost we
have not paid.

\subsection{What would settle the open questions}

The dependence induced by an estimated centre is the first. It is measured, its
form is classical, and a construction restoring exchangeability under it would
convert the ESS analysis from an approximation into an application of the
theorem. The second is the structure family: the diagnostic reported in
Section~\ref{sec:ess} is exploratory and its association with the payoff needs a
survey system not used to find it.
